% =====================================================================
%  Section 3 — Privacy Model, Channel-Importance Allocation, and
%               Honest Multi-Teacher Aggregation
%
%  Drop-in LaTeX for the privacy / methods section of the WACV paper.
%  No proprietary packages: needs only amsmath, amssymb, booktabs.
%
%  Numbers in Table 1 come from:
%      phase1_distillation/scripts/privacy_accounting.py
%  Numbers in Table 3 (PATE K-sweep) come from:
%      phase1_distillation/drive_pate_K_saturation_results.json
%      (5 seeds × 3 ε × 5 K values, paired student-init test).
% =====================================================================

\section{Privacy Model, Channel-Importance Allocation, and Honest Multi-Teacher Aggregation}
\label{sec:privacy}

We formalise the framework and the mechanism in four steps.
\S\ref{ssec:threat} states the threat model and explains why it differs
from the implicit \emph{per-iteration release} model used by prior DP
feature distillation work. \S\ref{ssec:accounting} gives the composition
accounting. \S\ref{ssec:wf} derives the channel-importance water-filling
(CANAL) allocation in closed form, bounds its advantage over uniform by
$R^{1/4}$ where $R$ is the importance spectrum spread, and notes that
this advantage is empirically negligible on RGB-fundus bottleneck
features ($R \!\approx\! 2$). \S\ref{ssec:pate} introduces the
\emph{honest multi-teacher PATE aggregation} mechanism that we propose
as the empirical contribution: training $K$ teachers on disjoint patient
cohorts and aggregating their bottlenecks before noising reduces the
per-channel $L_2$ sensitivity from $2$ to $2/K$, delivering a linear
$K$-fold reduction in the released noise scale.

% ---------------------------------------------------------------------
\subsection{Threat model: sample-once-per-image release}
\label{ssec:threat}

We consider a federated medical setting in which a data-holding site
processes each patient image \emph{once} through the teacher encoder,
adds Gaussian noise to the resulting bottleneck representation, and
releases the noisy representation. A separate party trains the student
against these fixed noisy targets. Formally, for training set
$\mathcal{D}=\{x_i\}_{i=1}^{N}$ and teacher encoder
$\phi_T : \mathcal{X}\!\to\!\mathbb{R}^{C\times H\times W}$, the
released objects are
\begin{equation}
\widetilde{z}_i \;=\; \mathrm{denorm}\!\bigl(\,\mathrm{norm}(\phi_T(x_i)) + \eta_i\,\bigr),
\qquad \eta_i \sim \mathcal{N}(0,\,\mathrm{diag}(\sigma^2)),
\label{eq:release}
\end{equation}
where $\mathrm{norm}(\cdot)$ performs per-channel $L_2$ clipping and
normalisation so that each channel of the released tensor lies in the
unit ball before noise is added.

We assume \textbf{per-patient privacy with at most one image per
patient} (DRIVE: $k\!=\!1$; BraTS: $k\!=\!1$). Both images and ground-
truth labels are treated as private; the student's training loop does
not access labels (Section~\ref{ssec:impl}), so the released features
are the only object the privacy mechanism must protect.

\textbf{Why this differs from prior work.}
Existing implementations of DP feature distillation re-sample noise
\emph{at every training iteration} of the student
\cite{lan2024gkd,Anonymous_anonymous}. Under this implicit
\emph{per-iteration} model, each training image's features are
released $Q = TB/N$ times across $T$ iterations of batch size $B$;
adaptive composition over $Q$ releases blows the user-level $\varepsilon$
up by orders of magnitude (cf.\ Table~\ref{tab:privacy}). The
sample-once model in \eqref{eq:release} eliminates this composition
entirely: each user contributes exactly one release.

% ---------------------------------------------------------------------
\subsection{Composition and budget accounting}
\label{ssec:accounting}

We use zero-Concentrated Differential Privacy (zCDP)
\cite{bun2016concentrated} as the working unit. For a Gaussian
mechanism on a $C$-dimensional vector with per-coordinate sensitivity
$\Delta_c$ and per-coordinate noise scale $\sigma_c$, the per-release
zCDP cost is
\begin{equation}
\rho_{\mathrm{rel}}
 \;=\; \sum_{c=1}^{C}\frac{\Delta_c^2}{2\,\sigma_c^2}.
\label{eq:rho-rel}
\end{equation}
After per-channel $L_2$ normalisation each channel has norm at most
$1$, so the replace-one $L_2$ sensitivity is the data-independent
constant $\Delta_c = 2/K$ for all $c$, where $K\geq 1$ is the
PATE-style teacher ensemble size (\S\ref{ssec:pate}).

\paragraph{Conversion to $(\varepsilon,\delta)$-DP.}
By \cite[Prop.~1.3]{bun2016concentrated}, a $\rho$-zCDP mechanism is
$(\rho + 2\sqrt{\rho\log(1/\delta)},\,\delta)$-DP for any $\delta>0$.
Conversely, the inverse map
\begin{equation}
\rho(\varepsilon,\delta)
 \;=\;\Bigl(\sqrt{\log(1/\delta)+\varepsilon}-\sqrt{\log(1/\delta)}\,\Bigr)^{2}
\label{eq:eps-to-rho}
\end{equation}
gives the maximum zCDP that fits inside a target
$(\varepsilon,\delta)$-DP budget.

\paragraph{Composition.}
The two threat models differ only in the number of times each user's
data participates in a release.
\begin{itemize}
\item \textbf{Sample-once per image} (ours). Each user contributes a
single release. The user-level cost is
$\rho_{\mathrm{user}} = \rho_{\mathrm{rel}}$, giving
$\varepsilon_{\mathrm{user}} = \varepsilon_{\mathrm{rel}}$. Caps and
importance are computed on a public proxy (\S\ref{ssec:impl}), so no
additional release accounting is required.
\item \textbf{Per-iteration release} (prior work). Each user
participates in $Q = TB/N$ releases of fresh noise on fresh teacher
queries. Naive zCDP composition gives
$\rho_{\mathrm{user}} = Q\,\rho_{\mathrm{rel}}$, mapped back to
$(\varepsilon_{\mathrm{user}}, \delta)$ via
\cite[Prop.~1.3]{bun2016concentrated}.
\end{itemize}

\paragraph{Reported vs.\ true privacy.}
Table~\ref{tab:privacy} contrasts the user-level $\varepsilon$ implied
by each threat model for the configurations used throughout this paper
($T{=}40{,}000$, $B{=}4$, $N{=}20$, $\delta{=}10^{-5}$). The
per-iteration model inflates the privacy cost by 400--1800$\times$
relative to what is reported by the noise calibration; the sample-once
model in \eqref{eq:release} faithfully reports user-level privacy by
construction.

\begin{table}[t]
\centering
\small
\caption{User-level $\varepsilon$ implied by each threat model for a
range of per-release $\varepsilon$ values. The per-iteration model
(prior work's implicit model) yields privacy guarantees that are
hundreds to thousands of times weaker than reported; the sample-once
model (ours) reports the actual user-level guarantee.
Parameters: $T{=}40{,}000$, $B{=}4$, $N{=}20$, $\delta{=}10^{-5}$.}
\label{tab:privacy}
\begin{tabular}{c c r r}
\toprule
per-release $\varepsilon$ & threat model & releases/user & user-level $\varepsilon$ \\
\midrule
\multirow{2}{*}{$2$}  & per-iteration & $8000$ & $812.1$ \\
                      & sample-once   & $1$    & $2.0$   \\
\addlinespace[2pt]
\multirow{2}{*}{$4$}  & per-iteration & $8000$ & $2712.4$ \\
                      & sample-once   & $1$    & $4.0$    \\
\addlinespace[2pt]
\multirow{2}{*}{$8$}  & per-iteration & $8000$ & $9014.8$ \\
                      & sample-once   & $1$    & $8.0$    \\
\addlinespace[2pt]
\multirow{2}{*}{$16$} & per-iteration & $8000$ & $28569.6$ \\
                      & sample-once   & $1$    & $16.0$    \\
\bottomrule
\end{tabular}
\end{table}

% ---------------------------------------------------------------------
\subsection{CANAL: channel-importance water-filling, and its $R^{1/4}$ ceiling}
\label{ssec:wf}

Given a fixed per-release budget $\rho_{\mathrm{rel}}$, the canonical
allocation across channels minimises an importance-weighted
reconstruction error
\begin{equation}
\min_{\sigma_1,\ldots,\sigma_C > 0}
\sum_{c=1}^C s_c\,\sigma_c^{2}
\quad\text{s.t.}\quad
\sum_{c=1}^C\frac{\Delta_c^{2}}{2\,\sigma_c^{2}} \le \rho_{\mathrm{rel}},
\label{eq:opt}
\end{equation}
where $s_c \ge 0$ is a per-channel importance score. We use
$s_c = \mathrm{mean}_{i,j} |\partial\mathcal{L}_{\mathrm{task}}/\partial \phi_{T,c}^{i,j}|$,
the GKD-style gradient magnitude \cite{lan2024gkd}, computed on a
public retinal proxy (\S\ref{ssec:impl}). The Lagrangian solution is
closed-form.

\begin{theorem}[Channel-WF / CANAL optimal allocation]
\label{thm:wf}
For positive importance scores $\{s_c\}$, positive sensitivities
$\{\Delta_c\}$ and budget $\rho_{\mathrm{rel}}>0$, the unique
minimiser of \eqref{eq:opt} is
\begin{equation}
\sigma^{\star}_c \;=\; \kappa \,\sqrt{\Delta_c}\, s_c^{-1/4},
\qquad
\kappa \;=\; \sqrt{\frac{1}{2\rho_{\mathrm{rel}}}
                          \sum_{c=1}^{C}\Delta_c\,\sqrt{s_c}}.
\label{eq:wf}
\end{equation}
\end{theorem}

The full proof is in Appendix~\ref{app:proof}. The advantage of CANAL
over uniform allocation is concentrated in the importance spread.

\begin{corollary}[CANAL advantage is bounded by $R^{1/4}$]
\label{cor:r14}
Let $R = s_{\max}/s_{\min}$ be the importance spectrum spread. The
ratio $\sigma_c^{\star}/\sigma_{c'}^{\star}$ between the least and most
important channels is exactly $R^{1/4}$. Consequently the optimal
CANAL objective $J_{\mathrm{WF}}^{\star}$ relative to the uniform
objective $J_{\mathrm{unif}}$ satisfies $J_{\mathrm{WF}}^{\star} \le
J_{\mathrm{unif}}$ with equality whenever the $s_c$ are constant
(Cauchy--Schwarz).
\end{corollary}

\paragraph{Empirical consequence on RGB fundus.}
On the bottleneck of a U-Net teacher trained on DRIVE, the gradient
importance spread is $R \!\approx\! 2$, so $R^{1/4} \!\approx\! 1.19$.
The corresponding $\sigma$ ratio between the most and least important
channels is empirically $1.13\times$ (Figure~\ref{fig:wf_ratio}). This
is below the Dice sensitivity of a trained segmentation decoder: in a
paired student-distillation experiment with $5$ seeds per cell we
measure CANAL $-$ uniform Dice differences of $+0.0003 \pm 0.0028$,
$-0.0003 \pm 0.0029$, $+0.0001 \pm 0.0028$ at $\varepsilon\!\in\!\{2,
8, 16\}$ respectively (all $|t| < 0.5\sigma$,
Section~\ref{sec:experiments}). The CANAL closed form is the right
allocation in the importance-weighted-distortion sense, but its
empirical room on RGB fundus is bounded by the flat importance
spectrum. On multi-modal MRI, where the modality-induced channel
spread $R$ is one to two orders of magnitude larger, we expect this
budget allocation to surface; we treat this as a future direction.

CANAL's empirical limit on this task motivates a complementary
mechanism that improves the privacy--utility frontier through
\emph{sensitivity reduction} rather than budget reallocation.

% ---------------------------------------------------------------------
\subsection{Honest multi-teacher PATE feature aggregation}
\label{ssec:pate}

Our proposed mechanism leverages a federated setting in which $K$
sites each hold a disjoint cohort of patients. Each site trains its
own teacher on its own patients; at release time, all $K$ teachers'
clipped and normalised bottlenecks are aggregated by mean before
Gaussian noise is added. This is the feature-space analogue of
the PATE classification mechanism of Papernot et al.
\cite{papernot2017semi,papernot2018scalable}, transported to dense
segmentation features rather than discrete labels.

\paragraph{Mechanism.}
Let $\{\phi_T^{(k)}\}_{k=1}^K$ be the $K$ teacher encoders, each
trained on the disjoint cohort $\mathcal{D}_k$, with
$\bigcup_k \mathcal{D}_k = \mathcal{D}$ and $\mathcal{D}_k \cap
\mathcal{D}_{k'} = \emptyset$ for $k \ne k'$. Let $\widehat{c}^{(k)}$
denote teacher $k$'s per-channel $L_2$ cap (computed on its public
proxy). For each image $x_i$, define
\begin{equation}
\widetilde{z}_i^{\mathrm{PATE}}
 \;=\; \mathrm{denorm}\!\Bigl(\,
   \tfrac{1}{K}\!\sum_{k=1}^K \mathrm{norm}\!\bigl(\phi_T^{(k)}(x_i);\widehat{c}^{(k)}\bigr)
   + \eta_i \,\Bigr),
\quad \eta_i \sim \mathcal{N}\!\bigl(0,\mathrm{diag}(\sigma^2)\bigr).
\label{eq:pate-release}
\end{equation}

\begin{theorem}[Aggregate sensitivity under PATE]
\label{thm:pate-sens}
Assume each teacher $\phi_T^{(k)}$ depends only on $\mathcal{D}_k$, the
cohorts are disjoint, and each teacher's bottleneck is per-channel
$L_2$-clipped to norm $\le 1$ after normalisation. Then a replace-one
change of a patient on the union $\mathcal{D}$ affects exactly one
teacher; the per-channel $L_2$ sensitivity of the aggregate in
\eqref{eq:pate-release} is
\begin{equation}
\Delta_c \;=\; 2/K \qquad \text{for all } c.
\label{eq:pate-delta}
\end{equation}
\end{theorem}

\begin{proof}[Proof sketch]
Replacing a patient in cohort $k$ leaves $\phi_T^{(k')}(x_i)$ unchanged
for all $k' \ne k$. After per-channel clip-and-normalise, teacher
$k$'s contribution to the aggregate lies in the unit $L_2$ ball, so its
replacement changes by at most $2$ per channel; the mean then changes
by at most $2/K$ per channel.
\end{proof}

\paragraph{Noise-scale consequence.}
Plugging $\Delta_c = 2/K$ into \eqref{eq:rho-rel} with the uniform
allocation gives
\begin{equation}
\sigma^{\mathrm{PATE}}(K)
 \;=\; \frac{2}{K}\sqrt{\frac{C}{2\rho_{\mathrm{rel}}}},
\label{eq:pate-sigma}
\end{equation}
so the noise scale shrinks linearly in $K$ at fixed budget. This is
\emph{mechanism-level} rather than allocation-level: any noise shape
(uniform, CANAL, threshold) inherits the same $K$-fold reduction when
paired with PATE aggregation. The CANAL closed-form of
Theorem~\ref{thm:wf} compounds with \eqref{eq:pate-sigma} but does not
replace it.

\paragraph{Empirical $K$-saturation.}
On DRIVE ($N\!=\!20$ training patients, $K$ teachers each on $\lceil
N/K \rceil$ patients), we sweep $K \in \{1, 3, 5, 8, 10\}$ at
$\varepsilon \in \{2, 8, 16\}$ with $5$ seeds per cell.
Table~\ref{tab:pate-ksweep} reports the resulting vessel Dice and the
paired student-init lift over the $K\!=\!1$ baseline. PATE delivers
monotonically increasing lift in $K$ up to $K\!=\!10$ on this dataset,
with no observed saturation: the paired lift over $K\!=\!1$ at
$\varepsilon\!=\!16$ grows from $+0.020$ at $K\!=\!3$ to $+0.052$ at
$K\!=\!10$ ($t = 15.3$), and the $K\!=\!10$ student Dice
($0.648$) exceeds the single-teacher clean Dice ($0.620$), an
ensemble-style regularisation effect.

\begin{table}[t]
\centering
\small
\caption{PATE multi-teacher feature aggregation on DRIVE. Paired
student-init test (5 seeds per cell); lift is mean $\pm$ SEM over
seeds. Monotonic in $K$ at every $\varepsilon$; no saturation observed
through $K\!=\!10$.}
\label{tab:pate-ksweep}
\begin{tabular}{r r r l l l}
\toprule
$K$ & $\Delta$ & $\sigma$ at $\varepsilon\!=\!16$
                                  & Dice at $\varepsilon\!=\!2$
                                                       & Dice at $\varepsilon\!=\!8$
                                                                              & Dice at $\varepsilon\!=\!16$ \\
\midrule
$1$  & $2.000$ & $8.64$  & $0.581$ (baseline)         & $0.594$ (baseline)            & $0.596$ (baseline) \\
$3$  & $0.667$ & $2.88$  & $+0.008\pm 0.004\,(1.7\sigma)$ & $+0.016\pm 0.004\,(3.6\sigma)$ & $+0.020\pm 0.003\,(6.0\sigma)$ \\
$5$  & $0.400$ & $1.73$  & $+0.012\pm 0.004\,(2.9\sigma)$ & $+0.025\pm 0.005\,(5.2\sigma)$ & $+0.031\pm 0.004\,(7.5\sigma)$ \\
$8$  & $0.250$ & $1.08$  & $+0.020\pm 0.004\,(5.0\sigma)$ & $+0.031\pm 0.007\,(4.7\sigma)$ & $+0.042\pm 0.005\,(8.1\sigma)$ \\
$10$ & $0.200$ & $0.86$  & $+0.026\pm 0.005\,(5.3\sigma)$ & $+0.034\pm 0.004\,(8.0\sigma)$ & $\mathbf{+0.052\pm 0.003\,(15.3\sigma)}$ \\
\bottomrule
\end{tabular}
\end{table}

\paragraph{Cohort-size caveat.}
At $K\!=\!10$ each teacher trains on only $2$ patients; the
monotonicity in Table~\ref{tab:pate-ksweep} suggests the ensemble
effect dominates the per-teacher under-training in this regime, but
extrapolating to higher $K$ would require more patients than DRIVE
provides. In a deployed federated setting, each ``teacher'' is one
hospital with hundreds of patients, and the cohort-size constraint
does not bind; we report DRIVE numbers as a controlled-setting
validation, with the understanding that the lift in
Table~\ref{tab:pate-ksweep} is a \emph{lower bound} on what a true
multi-hospital deployment would achieve.

% ---------------------------------------------------------------------
\subsection{Implementation}
\label{ssec:impl}

The sample-once threat model and PATE aggregation of
\eqref{eq:pate-release} are realised by a single precompute pass
before student training begins. Each of the $K$ teachers' encoders is
run once per training image, the $K$ bottlenecks are clipped,
normalised, and averaged, Gaussian noise calibrated to
\eqref{eq:pate-sigma} is added, the result is denormalised, and the
cache is keyed by image identifier. Student training then loads these
fixed targets at every iteration, eliminating both fresh teacher
queries and fresh noise samples from the inner loop. The student
trains a feature-matching loss against the cached target plus a
task-loss against the ground-truth mask; per-channel caps
$\widehat{c}^{(k)}$ and (for CANAL) importance scores
$\widehat{s}^{(k)}$ are computed once on the public proxy and
contribute zero additional privacy cost. Pseudocode appears in
Appendix~\ref{app:pseudocode}.
